% VI/45 reader bridge. Additive exposition; all original exercises retained.
\paragraph{What a blow-up remembers.}
A blow-up does not simply replace a point by an unrelated projective
space.  It records the limiting directions with which nearby points
approach the center.  Algebraically, for an ideal $I\subseteq A$,
the blow-up of $X=\operatorname{Spec}A$ along $I$ is
\[
 \operatorname{Bl}_I(X)=\operatorname{Proj}_A
 \bigl(A\oplus I T\oplus I^2T^2\oplus\cdots\bigr).
\]
The graded algebra is the \emph{Rees algebra}; retaining the ideal
powers matters because they encode the scheme structure of the center.
The blow-up is an isomorphism over the complement of $V(I)$.  Its
inverse image of $V(I)$ is the exceptional locus, whose geometry
requires a local calculation rather than a set-theoretic guess.

\paragraph{Worked example 1: blowing up the affine plane at the origin.}
Put $A=k[x,y]$ and $I=(x,y)$.  Using projective direction
coordinates $[U:V]$ gives the closed subscheme
\[
 B=\{((x,y),[U:V])\in\mathbb A_k^2\times\mathbb P_k^1:
               xV-yU=0\}.
\]
The equation says that $[U:V]=[x:y]$ away from the origin.
On $U\ne0$, write $t=V/U$.  The relation becomes $y=xt$, so
\[
 B\cap D_+(U)=\operatorname{Spec}k[x,t],\qquad
 (x,t)\longmapsto(x,xt).
\]
On $V\ne0$, write $s=U/V$.  Here $x=ys$, and
\[
 B\cap D_+(V)=\operatorname{Spec}k[s,y],\qquad
 (s,y)\longmapsto(ys,y).
\]
On the overlap, $s=t^{-1}$ and $y=xt$.  These charts are
ordinary affine planes glued by explicit coordinate changes.
Over $(x,y)=(0,0)$ the projective coordinate $[U:V]$ is free;
hence the exceptional fiber is $E\simeq\mathbb P^1_k$.
The total space is smooth, but it maps an entire curve to one point.
The normal bundle of $E$ in $B$ is $\mathcal O_{\mathbb P^1}(-1)$:
the local equations $x$ and $y$ obey $y=tx$ on the overlap,
which is the characteristic transition for the conormal/normal
pair (with the normal line bundle of degree $-1$).

\paragraph{Worked example 2: strict transform of a cusp.}
Let $C\subseteq\mathbb A_k^2$ be defined by
$y^2=x^3$; for simplicity let $\operatorname{char}k\ne2$.
Substitute $y=xt$ on the first blow-up chart:
\[
 y^2-x^3=x^2(t^2-x).
\]
The total transform contains the exceptional factor $x^2$.
The \emph{strict transform} is instead $C':t^2=x$, obtained by
closing the inverse image of $C\setminus\{0\}$; it is smooth,
because its defining equation has derivative $-1$ with respect
to $x$.  Its intersection with $E:x=0$ is the double point
$\operatorname{Spec}k[t]/(t^2)$.  Thus one blow-up makes the
strict transform smooth but does \emph{not} produce a transverse
intersection with the exceptional curve.  Resolving the curve's
singularity and producing a normal-crossings embedded resolution
are different requirements.

\paragraph{Worked example 3: a node remembers two tangent directions.}
For $C:y^2=x^2(x+1)$, the same chart gives
\[
 x^2\bigl(t^2-(x+1)\bigr)=0,
 \qquad C':\ t^2=x+1.
\]
Assume $\operatorname{char}k\ne2$.  At $x=0$ the strict transform
meets the exceptional curve in the two distinct points $t=1$ and
$t=-1$.  The two tangent directions of the node have become two
separate points, in contrast with the double intersection of the cusp.
This also explains why one should divide out the appropriate
exceptional factor before assessing the strict transform.

\paragraph{Universal property and the next chapter.}
The blow-up is universal among morphisms $f:Y\to X$ for which
the extended ideal $I\mathcal O_Y$ is an invertible ideal sheaf
(with the usual interpretation when the pullback ideal is
locally principal and generated by a non-zerodivisor).  This
property ties together the Rees-algebra and affine-chart pictures.
The exceptional curve explains the geometry of a birational map;
VI/46 turns instead to the extension of sections of sheaves.
